\section{Mathematics: everyday formulas}\label{math:basics}
\subsection{Sums, ranges and averages}
\hint{$n\ge0$ is an integer; $a$ is the first term, $d$ the step,
$r$ the ratio. Empty sums are 0.}
\[
\begin{aligned}
 \sum_{i=1}^{n}i&=\frac{n(n+1)}2,\\
 \sum_{i=1}^{n}i^2&=\frac{n(n+1)(2n+1)}6,\\
 \sum_{i=1}^{n}i^3&=\left(\frac{n(n+1)}2\right)^2,\\
 \sum_{i=0}^{n-1}(a+id)&=na+\frac{dn(n-1)}2.
\end{aligned}
\]
\[
 \sum_{i=0}^{n-1}r^i=\frac{r^n-1}{r-1}\quad(r\ne1);
 \qquad r=1:\ n.
\]
For $|r|<1$, $\sum_{i=0}^{\infty}r^i=1/(1-r)$.
Modulo an integer, division requires an invertible denominator
or an exact-division method; do not divide reduced residues normally.

For integers $L\le R$ and $k>0$, multiples of $k$ in $[L,R]$:
\[
 \left\lfloor R/k\right\rfloor-
 \left\lfloor(L-1)/k\right\rfloor.
\]
This also works for negative endpoints; use floor division.
For real values $a_1,\ldots,a_n$, $n\ge1$, minimizing over real $x$:
$\sum_i|a_i-x|$ is minimized by a median;
$\sum_i(a_i-x)^2$ by the mean $x=\sum_i a_i/n$.

\subsection{Counting choices}\label{math:counting}
\hint{Objects are distinct unless stated. $n,k,s,u$ are nonnegative
integers. $0!=1$; $\binom nk=n!/[k!(n-k)!]$ for $0\le k\le n$,
and 0 outside that range or when the upper argument is negative.}
\begin{tabularx}{\linewidth}{@{}Xl@{}}
\toprule
What is counted & Count\\
\midrule
Orders of $n$ objects & $n!$\\
Choose $k$, order matters ($k\le n$) & $n!/(n-k)!$\\
Choose $k$, order ignored & $\binom nk$\\
Subsets of $n$ objects & $2^n$\\
Length-$k$ words over $n$ symbols & $n^k$\\
Paths: $a$ right, $b$ up & $\binom{a+b}{a}$\\
Choose $k$ of $n\ge1$ types, repeats & $\binom{n+k-1}{k}$\\
Circular orders, rotations identified ($n\ge1$) & $(n-1)!$\\
\bottomrule
\end{tabularx}
The empty word has one realization, including for an empty alphabet:
use $0^0=1$ in these counting formulas. Circular orders here distinguish
reflections; necklaces with repeated colors need Burnside's lemma.

\textbf{Repeated labels.} If type $i$ occurs $c_i\ge0$ times and
$\sum_i c_i=n$, distinct orders number $n!/\prod_i c_i!$.

\textbf{Stars and bars.} For $k\ge1$ labeled boxes and total $s$:
\[
\begin{aligned}
 \#\{x_1+\cdots+x_k=s,\ x_i\ge0\}
   &=\binom{s+k-1}{k-1},\\
 \#\{x_1+\cdots+x_k=s,\ x_i\ge1\}
   &=\binom{s-1}{k-1}\quad(s\ge k).
\end{aligned}
\]
The second count is 0 for $s<k$. With $0\le x_i\le u$,
\[
 \sum_{j=0}^{k}(-1)^j\binom kj
      \binom{s-j(u+1)+k-1}{k-1}.
\]

\subsection{Catalan, derangements and coverage}
\textbf{Catalan $C_n$.} Balanced parentheses with $n$ pairs,
full ordered binary trees with $n$ internal nodes, or triangulations
of a convex $(n+2)$-gon ($n\ge1$):
\[
 C_n=\frac{1}{n+1}\binom{2n}{n}\quad(n\ge0).
\]
$C_0,C_1,\ldots=1,1,2,5,14,42,132,429,\ldots$.

\textbf{Derangements $D_n$.} Permutations of $n$ labels with no label
in its original position:
\[
 D_0=1,\quad D_1=0,\quad
 D_n=(n-1)(D_{n-1}+D_{n-2})\quad(n\ge2).
\]
$D_0,D_1,\ldots=1,0,1,2,9,44,265,\ldots$.

\textbf{Coupon collector.} Independent uniform draws from $n\ge1$
types, with replacement: expected draws to see all types is
$nH_n$, where $H_n=\sum_{i=1}^n1/i$.

\subsection{Vectors, angles and distances}\label{math:geometry}
\hint{Points/vectors are in the plane. Write $a=(a_x,a_y)$,
$b=(b_x,b_y)$; $|a|$ is Euclidean length.}
\[
 a\cdot b=a_xb_x+a_yb_y,\qquad
 a\times b=a_xb_y-a_yb_x.
\]
$a\times b>0$: $b$ is counterclockwise from $a$ through an angle
less than $\pi$; zero means parallel. Orientation of three points
$p,q,r$ is the sign of $(q-p)\times(r-p)$.
$a\cdot b=0$ means perpendicular for nonzero vectors.
\[
 |p-q|^2=(p_x-q_x)^2+(p_y-q_y)^2.
\]
Compare squared distances when possible. Integer coordinates allow
exact predicates; bound intermediate products and accumulated sums.
Python integers are exact; C++ may need a wider intermediate type.

For nonzero vectors, the unsigned angle between them is
$\operatorname{atan2}(|a\times b|,a\cdot b)\in[0,\pi]$.
Direction of $a\ne0$ is $\operatorname{atan2}(a_y,a_x)$.
Convert degrees to radians by multiplying by $\pi/180$.
Counterclockwise rotation through $\theta$ radians:
\[
 (x,y)\mapsto(x\cos\theta-y\sin\theta,
              x\sin\theta+y\cos\theta).
\]

\textbf{Point $p$ to line/segment $ab$.} Require $a\ne b$;
put $v=b-a$, $t=(p-a)\cdot v/(v\cdot v)$.
For the line, projection is $q=a+tv$.
For the segment, first clamp $t$ to $[0,1]$.
Distance is $|p-q|$. If $a=b$, use $|p-a|$.

\subsection{Polygon area and lattice points}
\hint{$p_0,\ldots,p_{n-1}$ are vertices in boundary order,
$p_n=p_0$; $A$ denotes nonnegative area.}
\[
 A_2=\sum_{i=0}^{n-1}p_i\times p_{i+1},\qquad A=|A_2|/2.
\]
For a simple polygon, $A_2>0$ means counterclockwise orientation.
Triangle area is $|(b-a)\times(c-a)|/2$.
A simple $n$-vertex polygon has interior-angle sum $(n-2)\pi$.

\textbf{Pick's theorem.} For a simple polygon with integer vertices
and no holes, let $I$ be its strictly interior lattice-point count and
$B$ its boundary lattice-point count:
\[
 A=I+B/2-1,\qquad
 B=\sum_{\text{edges}}\gcd(|\Delta x|,|\Delta y|).
\]
Thus $I=(|A_2|-B+2)/2$. The differences are between the endpoints
of each edge; do not count shared vertices separately.

\subsection{Triangles, circles and spheres}
\hint{A nondegenerate triangle has side lengths $a,b,c>0$,
opposite angles $\alpha,\beta,\gamma$, semiperimeter
$s=(a+b+c)/2$ and area $A>0$.}
\[
\begin{aligned}
 c^2&=a^2+b^2-2ab\cos\gamma,\\
 A&=\sqrt{s(s-a)(s-b)(s-c)},\\
 R&=abc/(4A),\qquad r=A/s,\\
 a/\sin\alpha&=b/\sin\beta=c/\sin\gamma=2R.
\end{aligned}
\]
$R$ is circumradius; $r$ is inradius. With $a\le b\le c$,
existence requires $a+b>c$; the largest angle is right, acute or
obtuse according as $a^2+b^2-c^2$ is zero, positive or negative.

For a circle of radius $r\ge0$ and central angle $0\le\theta\le2\pi$
in radians: circumference $2\pi r$, area $\pi r^2$, arc length $r\theta$,
sector area $r^2\theta/2$, chord length $2r\sin(\theta/2)$.
A sphere of radius $r$ has surface area $4\pi r^2$ and volume
$4\pi r^3/3$.

\hint{For floating-point geometry, choose tolerances for the input
scale; one fixed epsilon is not universal. Clamp a cosine to
$[-1,1]$ before \texttt{acos}, and a tiny negative roundoff to 0
before \texttt{sqrt}. Never use an epsilon comparator for sorting
unless it still defines a strict weak ordering.}
